%% title: English Subject-Verb Agreement
%% category: phenomena

English verbs agree with their subjects in person and number. I assume like most that `subject' here refers to whatever DP sits in specTP, and that it is T which agrees with the subject:

\pex \dref{eng-sva-lemurs}
  \a Lemurs like lemonade.
  \a \begin{forest}
[TP
  [T\feat{+pl}]
  [VoiceP
    [DP\feat{+pl} [\textit{lemurs}, roof]]
    [Voice\1
      [Voice]
      [\vP
        [DP\feat{–pl} [\textit{lemonade}, roof]]
        [\lv\1
          [\lv]
          [\rt{LIKE}]]]]]]
\end{forest}
\xe

In the labelling framework, the analysis is simple. First the lexical root \rt{LIKE} merges with a verbalizer, and the resulting mother node is unvalued for number, because its children are both unvalued for number (this is obligatory if the algorithm is \textbf{conservative}).

Then, the resulting node is merged with the object DP \textit{lemonade}, which yields a mother node which has the [–pl] feature, because \textit{lemonade} has a [–pl] feature:

\pex \dref{eng-sva-vp}
  \a \begin{forest}
[\vP\feat{–pl}
  [DP\feat{–pl} [\textit{lemonade}, roof]]
  [\lv\1
    [\lv]
    [\rt{LIKE}]]]
\end{forest}
  \a $\tau(\text{[–pl]}, \text{[ ]}) = \text{[–pl]}$
\xe

Next, the resulting \vP merges with Voice, and continues to project its [–pl] feature, because Voice has no number feature intrinsically.

Then, you merge the subject, \textit{lemurs}, which has a [+pl] feature. Here, there must be a rule telling this merge combination to project [+pl] and not [–pl]:

\pex \dref{eng-sva-voicep}
  \a \begin{forest}
[VoiceP\feat{+pl}
  [DP\feat{+pl} [\textit{lemurs}, roof]]
  [Voice\1\feat{–pl}
    [Voice]
    [\vP\feat{–pl} [\textit{like lemonade}, roof]]]]
\end{forest}
  \a $\tau(\text{[+pl]}, \text{[–pl]}) = \text{[+pl]}$
\xe

Finally, you merge T, and the mother of T and VoiceP also receives [+pl], as T has no intrinsic number feature.

\pex \dref{eng-sva-tp}
  \a \begin{forest}
[TP\feat{+pl}
  [T]
  [VoiceP\feat{+pl} [\textit{lemurs like lemonade}, roof]]]
\end{forest}
  \a $\tau(\text{[ ]}, \text{[+pl]}) = \text{[+pl]}$
\xe

But it is T which needs to receive the [+pl] feature in the end. The simplest way to deal with this is to say that TP and T are really the same object -- we are implicitly using Bare Phrase Structure -- and so the features a head is eventually pronounced with are really just the features of its maximal projection. I'll call this \textbf{maximal-projection sharing}, or \textbf{MP-sharing}, for now.

There are other ways to analyze this phenomenon within the labelling framework, but this is the simplest.
